\subsection{Identify the project authority, sponsor/client, project manager, decision responsibilities, and escalation mechanism.}

\subsubsection{Client, Project Authority và Project Sponsor}

Nhà trường là \textit{Client} của dự án, với nhu cầu xây dựng một nền tảng tập trung để quản lý sự kiện, đăng ký tham gia, điểm danh và tổng hợp báo cáo. Theo phản hồi trong quá trình tham vấn, thầy Nguyễn Minh Tâm là đại diện \textit{Client} trong phạm vi bài tập lớn, chịu trách nhiệm xác nhận yêu cầu, phạm vi, tiêu chí nghiệm thu và các thay đổi lớn.

Để xác định cơ cấu quản lý dự án, nhóm giả định nhà trường đồng thời là \textit{Project Sponsor} và đề xuất thầy Nguyễn Minh Tâm đảm nhiệm vai trò \textit{Project Authority}. Trong cơ cấu này, thầy xem xét và xác nhận \textit{Project Baseline}, quyết định các vấn đề vượt quá quyền hạn của \textit{Project Manager}, đồng thời đại diện \textit{Project Sponsor} xem xét dự toán và nhu cầu nguồn lực do nhóm đề xuất. Ngân sách trong kế hoạch phục vụ việc đánh giá và quản lý bài tập lớn, không thể hiện cam kết cấp kinh phí thực tế của nhà trường.

\subsubsection{Project Manager}

Nguyễn Đăng Khoa là \textit{Project Manager}, chịu trách nhiệm tổ chức và điều phối nhóm trong suốt quá trình thực hiện dự án. Các trách nhiệm chính gồm lập kế hoạch, phân công công việc, điều phối nguồn lực, theo dõi tiến độ và công sức thực tế, quản lý rủi ro và vấn đề phát sinh, duy trì trao đổi với đại diện \textit{Client}, đồng thời lập dự toán và tổng hợp các vấn đề cần \textit{Project Authority} quyết định.

Trước khi trình xác nhận \textit{Project Baseline}, \textit{Project Manager} kiểm tra tính nhất quán giữa phạm vi, sản phẩm bàn giao, hoạt động, ước tính, lịch trình và tiêu chí nghiệm thu. Trong quá trình triển khai, \textit{Project Manager} so sánh kết quả thực tế với kế hoạch, xác định nguyên nhân sai lệch và tổ chức điều chỉnh trong phạm vi quyền hạn được giao.

Ngoài trách nhiệm quản lý, Nguyễn Đăng Khoa kiêm hỗ trợ một số tác vụ Backend theo phân công tại Mục~1.7.1; thời gian này được tính chung vào tải công việc của thành viên PM và không thay đổi thẩm quyền phê duyệt, nghiệm thu đã xác định.

\subsubsection{Trách nhiệm ra quyết định}

Trách nhiệm ra quyết định được phân theo ba cấp:

\begin{itemize}
    \item \textbf{Cấp nhóm:} các thành viên đưa ra quyết định kỹ thuật dưới sự điều phối của \textit{Project Manager}, căn cứ vào yêu cầu chức năng, yêu cầu chất lượng và nguồn lực hiện có.

    \item \textbf{Cấp Project Manager:} Nguyễn Đăng Khoa quyết định việc điều hành hằng ngày, điều chỉnh phân công, phối hợp nhân sự và thứ tự thực hiện hoạt động, miễn là các quyết định vẫn bảo đảm những cam kết trong \textit{Project Baseline}.

    \item \textbf{Cấp Project Authority:} thầy Nguyễn Minh Tâm xem xét và xác nhận các thay đổi bổ sung hoặc loại bỏ chức năng đã cam kết, điều chỉnh tiêu chí nghiệm thu, thay đổi mốc bàn giao hoặc những thay đổi có ảnh hưởng đáng kể đến chi phí và nguồn lực.
\end{itemize}

\subsubsection{Cơ chế escalation}

Khi phát hiện vấn đề, thành viên phụ trách ghi nhận nội dung, ảnh hưởng dự kiến và biện pháp xử lý ban đầu, sau đó báo cáo cho \textit{Project Manager}. \textit{Project Manager} đánh giá ảnh hưởng của vấn đề đến phạm vi, tiến độ, công sức, chi phí và chất lượng. Các vấn đề thuộc phạm vi điều hành nội bộ được giải quyết bằng cách điều chỉnh công việc và phối hợp thành viên.

Nếu vấn đề vượt quá quyền hạn, không thể xử lý bằng nguồn lực hiện có hoặc có nguy cơ ảnh hưởng đến cam kết với \textit{Client}, \textit{Project Manager} thực hiện \textit{escalation} lên thầy Nguyễn Minh Tâm. Báo cáo \textit{escalation} cần nêu rõ tình trạng, nguyên nhân, mức độ ảnh hưởng, các phương án khả thi và khuyến nghị của nhóm. Sau khi có quyết định, \textit{Project Manager} cập nhật kế hoạch và điều phối nhóm thực hiện.

\subsubsection{Ghi nhận quyết định và quản lý Project Baseline}

Các quyết định quan trọng được ghi trong nhật ký dự án hoặc biên bản trao đổi với khách hàng, bao gồm người ra quyết định, nội dung được chấp thuận và trách nhiệm thực hiện. Nếu quyết định làm thay đổi \textit{Project Baseline}, nhóm cập nhật phạm vi, sản phẩm bàn giao, hoạt động, nguồn lực, chi phí và lịch trình liên quan, sau đó phát hành phiên bản \textit{baseline} mới khi được thầy Nguyễn Minh Tâm xác nhận.

Các phiên bản trước được lưu lại để đối chiếu với kế hoạch điều chỉnh và kết quả thực tế, qua đó bảo đảm khả năng truy vết thay đổi.